% ============================================================
% CAPÍTULO 7 — Propuesta de Cierre
% ============================================================
\chapter{Propuesta de Cierre}
\label{cap:propuesta}

A partir de los an\'alisis previos, este cap\'itulo presenta cuatro
opciones de cierre contempladas, la opci\'on recomendada por el equipo
y un plan de entrega concreto de los artefactos del proyecto.

\section{Opciones contempladas}

\subsection*{Opci\'on A --- Archivo limpio}

Cierre completo y formal del proyecto, con publicaci\'on del repositorio
en GitHub bajo licencia abierta y entrega de la documentaci\'on como
material p\'ublico. Cero coste de mantenimiento posterior.

\textbf{Ventajas}: simple, definitivo, sin compromisos pendientes.
\textbf{Limitaciones}: no preserva opciones expl\'icitas de continuidad.

\subsection*{Opci\'on B --- Cierre con opci\'on de reactivaci\'on}

Como Opci\'on A, m\'as un anexo que detalla qu\'e condiciones
permitir\'ian una reactivaci\'on parcial del proyecto (financiaci\'on
m\'inima requerida, equipo m\'inimo necesario, alcance reducido viable).

\textbf{Ventajas}: mantiene puerta abierta sin compromiso operativo.
\textbf{Limitaciones}: si las condiciones nunca se dan, equivale a A.

\subsection*{Opci\'on C --- Transferencia a otro grupo}

Cesi\'on formal del proyecto a un grupo de investigaci\'on o a otro
equipo capaz de mantenerlo. Requiere acuerdo expl\'icito de
transferencia, atribuci\'on y, eventualmente, soporte de los autores
durante el traspaso.

\textbf{Ventajas}: el sistema sigue vivo bajo otra estructura.
\textbf{Limitaciones}: requiere existencia de receptor interesado y
mecanismos formales de transferencia.

\subsection*{Opci\'on D --- Reducci\'on a caso docente}

Conversi\'on del proyecto en material did\'actico estable: el sistema
se documenta como caso de estudio reutilizable para asignaturas de IoT,
sistemas embebidos, agronom\'ia digital y gesti\'on de proyectos. No
requiere mantenimiento operacional ni infraestructura activa.

\textbf{Ventajas}: encaja perfectamente con el marco de Proyecto
Docente original; cumple objetivos did\'acticos; cero coste recurrente.
\textbf{Limitaciones}: el sistema deja de estar desplegado
operativamente.

\section{Opci\'on recomendada: combinaci\'on A + D}
\label{sec:recomendacion}

El equipo propone articular el cierre como \textbf{combinaci\'on de
archivo limpio (Opci\'on A) y reducci\'on a caso docente (Opci\'on D)},
por las siguientes razones:

\begin{itemize}[leftmargin=*]
    \item Encaja con el marco original de Proyecto Docente bajo el que
    se obtuvo la financiaci\'on inicial.
    \item Cumple los objetivos did\'acticos planteados (desarrollo de
    competencias multidisciplinares, generaci\'on de material
    aprovechable docente).
    \item No requiere mantenimiento operacional ni financiaci\'on
    recurrente.
    \item Preserva el valor acad\'emico generado.
    \item Mantiene abierta la posibilidad de retomar componentes
    espec\'ificos como TFG/TFM futuros sobre piezas concretas
    (firmware, SDK, cl\'iuster cloud, etc.).
\end{itemize}

\subsection{Concreci\'on de la propuesta}

\begin{enumerate}[leftmargin=*]
    \item \textbf{Repositorio p\'ublico}: publicaci\'on completa del
    c\'odigo en GitHub bajo licencia \textbf{MIT} para el
    c\'odigo y \textbf{CC-BY-SA 4.0} para la documentaci\'on, con
    atribuci\'on a la Asociaci\'on ESIBot y a la Universidad de Sevilla.

    \item \textbf{Documentaci\'on como material reutilizable}: la
    \textit{Documentaci\'on T\'ecnica del Sistema} (84~p\'aginas) y el
    presente informe quedan disponibles, con licencia abierta, para
    cualquier uso docente que la Escuela o el Departamento consideren
    oportuno. No se propone aqu\'i ninguna asignatura concreta: la
    decisi\'on sobre d\'onde y c\'omo aprovechar el material corresponde a
    quien conozca los programas vigentes.

    \item \textbf{Disponibilidad puntual de los autores}: el equipo se
    compromete a sesiones puntuales de presentaci\'on del sistema a
    futuros estudiantes o equipos interesados, durante un periodo de
    \textbf{doce meses} tras el cierre formal.

    \item \textbf{Cat\'alogo de TFG/TFM derivables}: se propone elaborar
    un documento breve (1--2 p\'aginas) listando posibles trabajos
    derivables sobre componentes espec\'ificos del sistema, para uso
    de coordinadores docentes interesados.
\end{enumerate}

\section{Plan de entrega de artefactos}

El cronograma de entrega propuesto se articula en cuatro semanas tras
la aceptaci\'on formal del cierre, con los detalles t\'ecnicos
recogidos en el Anexo~\ref{anx:plan_entrega}.

\begin{table}[H]
    \centering
    \caption{Cronograma resumido de cierre.}
    \label{tab:cronograma}
    \begin{tabular}{ll}
        \toprule
        \textbf{Semana} & \textbf{Hito} \\
        \midrule
        Semana 1 & Publicaci\'on del repositorio bajo licencia abierta. \\
        Semana 2 & Entrega formal de documentaci\'on (PDF + fuentes LaTeX). \\
        Semana 3 & Devoluci\'on / cesi\'on del hardware adquirido. \\
        Semana 4 & Cierre de la infraestructura cloud (VM, dominio, IP). \\
        \bottomrule
    \end{tabular}
\end{table}

El equipo solicita, como reconocimiento simb\'olico del cierre, un
\textbf{acta de cierre formal} firmada por la Universidad y la
Asociaci\'on, que recoja la conclusi\'on del Proyecto Docente y los
artefactos entregados. Este documento constituye reconocimiento del
trabajo realizado y queda como referencia institucional de la
colaboraci\'on.
